% =========================================================
% SECCIÓN 2.7 - CONCLUSIONES DEL CAPÍTULO
% =========================================================

En este capítulo se estableció la base instrumental y metrológica sobre la cual descansa el sistema híbrido de predicción de emisiones. Se diseñó, implementó y caracterizó un módulo de adquisición basado en sensores NDIR de bajo costo para \ce{CO2} y un sensor semiconductor para \ce{CH4}, integrado en un nodo IoT con procesamiento en el borde mediante ESP32. El protocolo de ventanas cerradas de \SI{30}{\minute}, con intercalación de fases de aireación, permitió cuantificar tasas de emisión sin perturbar la oxigenación del reactor ni inducir hipoxia larval, mientras que la derivación de conversión ppm$\rightarrow$g/min habilitó la expresión de flujos másicos compatibles con el balance de carbono del modelo bioenergético.

La calibración frente a patrones certificados y la compensación cruzada por temperatura y humedad relativa redujeron el MAPE del sensor NDIR del \SI{18.4}{\percent} (crudo) al \SI{9.8}{\percent} (corregido), con una correlación de Pearson $R = 0{,}997$ respecto al patrón y una incertidumbre expandida del \SI{4.9}{\percent} para concentraciones de \ce{CO2}. Estos valores satisfacen holgadamente los umbrales de la hipótesis H1 ($R > 0{,}85$, MAPE $<$ \SI{15}{\percent}) y validan la viabilidad de sensores de bajo costo para monitoreo continuo en entornos agroindustriales tropicales, donde la infraestructura de referencia tradicional resulta inaccesible \parencite{biagiDevelopmentMachineLearningbased2024, rajRealTimeEstimation2023}.

Se identificaron, asimismo, las limitaciones estructurales del hardware que condicionan el alcance del sistema. El canal de \ce{CH4}, afectado por sensibilidad cruzada con humedad y VOCs, no alcanzó resolución suficiente para cuantificar flujos másicos continuos; por ello, se operó como clasificador binario de eventos de anaerobiosis (umbral \SI{100}{ppm}), función que, si bien sacrifica la cuantificación, preserva la utilidad diagnóstica operativa prevista para el Capítulo 5. La deriva en condiciones de HR $>$ \SI{90}{\percent} y la no-linealidad residual por encima de \SI{4000}{ppm} quedan documentadas como fuentes de incertidumbre a monitorear durante la validación experimental de los capítulos siguientes \parencite{herrinCollateralDataQuality2021, rossiEstimatingDynamicsGreenhouse2024}.

La trazabilidad del dato crudo quedó garantizada mediante la cadena documentada: lectura cruda $\rightarrow$ corrección de offset $\rightarrow$ compensación cruzada ($T$, HR) $\rightarrow$ conversión a flujo másico con incertidumbre propagada. Esta trazabilidad habilita la integración del Capítulo 2 con el modelo DEB del Capítulo 3 ---donde los flujos $\dot{m}_{\ce{CO2}}(t)$ alimentarán la validación de las ecuaciones de estado--- y con el estimador híbrido DEB--GP del Capítulo 4, donde la incertidumbre instrumental se combinará con la varianza del proceso gaussiano para producir intervalos de predicción robustos \parencite{eriksenDynamicModellingFeed2022, padmanabhaComprehensiveDynamicGrowth2020}.

En síntesis, el capítulo entrega un flujo de datos crudos validados, con error cuantificado y trazable, que constituye la entrada indispensable para cerrar el ciclo del sistema híbrido: sensores $\rightarrow$ modelo $\rightarrow$ estimador $\rightarrow$ métricas ambientales.
